\section{Introduction}
\label{sec:introduction}

Let $M^n$ be a closed connected oriented smooth manifold, let
$\Gamma=\pi_1(M)$, and let $c_M:M\to B\Gamma$ be a classifying map
for its universal cover. Here $B\Gamma$ is a CW $K(\Gamma,1)$: a connected
CW complex with fundamental group $\Gamma$ and contractible universal cover.
The manifold is
\emph{rationally essential} if
\[
 (c_M)_*[M]\ne0\quad\text{in }H_n(B\Gamma;\Q),
\]
and \emph{rationally inessential} otherwise. The fundamental class in this
definition has rational coefficients. The condition is independent of the
choice of classifying map. It records whether the fundamental class remains
visible after passing from the manifold to the topology of its fundamental
group.

We prove the rational-essentiality obstruction to positive scalar curvature
formulated in \cite[Conjecture~2.9]{Hanke2011}.

\begin{theorem}[Rational inessentiality]
\label{thm:main}
Let $M^n$ be a closed connected oriented smooth manifold of dimension
$n\ge2$. If $M$ admits a Riemannian metric of strictly positive scalar
curvature, then
\[
 (c_M)_*[M]=0\qquad\text{in }H_n(B\pi_1(M);\Q).
\]
\end{theorem}

The theorem has no spin, universal-cover spin, or fundamental-group
hypothesis, and no upper dimension bound. In particular, the group may
have torsion and its classifying space need not have a finite-dimensional
model.

A manifold is \emph{aspherical} if its universal cover is contractible.

\subsection{From a detected class to contracting maps}

In the later geometric convention, a closed oriented Riemannian
$n$-manifold is \emph{enlargeable} if, for every $\varepsilon>0$, some cover
with the lifted metric admits an $\varepsilon$-Lipschitz map of nonzero degree
to the unit sphere $S^n$ that is constant outside a compact set. This
definition imposes no spin condition. Rational essentiality supplies a
group-cohomology class whose pullback pairs nontrivially with $[M]$, but it
does not guarantee these contracting maps. Brunnbauer and Hanke construct
rationally essential manifolds which are not enlargeable, as well as
enlargeable manifolds whose universal covers alone do not admit such maps
at every scale, that is, are not hyperspherical
\cite[Theorems~1.4--1.5]{BrunnbauerHanke2010}. Their examples are not aspherical.

Guth and Braun--Sauer give a different comparison between topology and
metric largeness: at every radius, they guarantee a ball in the universal
cover whose volume is at least a dimensional multiple of the radius to the
$n$th power, for closed aspherical and rationally essential manifolds,
respectively \cite{Guth2011,BraunSauer2021}. These bounds do not have the
sharp Euclidean leading constant, so they do not determine the scalar-curvature
term in the small-ball volume expansion.

Our proof converts the cohomology pairing into a local degree by adding a
parameter manifold, then changes the metric to contract the resulting map
while retaining positive scalar curvature. The main estimate makes the scalar
cost grow only logarithmically with the number of comparison coordinates.
Suppose for contradiction that $\Scal_M\ge\kappa>0$ and that $M$ is
rationally essential. The construction has four stages.

\paragraph{Finite comparison and local degree.}
At each sufficiently large finite word-metric comparison scale $R$, an
auxiliary closed oriented manifold $C$ with a principal left $\Gamma$-cover
$\widehat C\to C$ gives the cover
$X=(\widetilde M\times\widehat C)/\Gamma\to M\times C$. The detecting cocycle
yields a compactly supported cohomology class on $X$ that evaluates
nontrivially on its locally finite fundamental class. Distance profiles supply a preliminary
map $\mathcal V:X\to\R^q$, for a finite integer
$q>n$. Together with the parameter projection $\bar o:X\to C$, it defines
the map $x\mapsto(\bar o(x),\mathcal V(x))$ into $C\times\R^q$.

The construction transfers the compact pairing to a relative cohomology class
on a compact region of this product, evaluated nontrivially on a relative
cycle carried by the preceding map. Relative duality gives a complementary
relative rational homology class, and rational realization represents a nonzero multiple
by a compact oriented $(q-n)$-manifold with boundary mapped to
$C\times\R^q$. Its transverse intersections with the preceding map have
nonzero signed count. Doubling the parameter manifold and reflecting its map
gives a closed oriented manifold $A$ with a map $(o_0,y_0):A\to C\times\R^q$.
Pulling back $\widehat C\to C$ along $o_0$ gives a principal cover
$\widehat A\to A$ and the fibre product
\[
 Z=X\times_C A=(\widetilde M\times\widehat A)/\Gamma
       \longrightarrow M\times A.
\]
This cover has dimension $q=n+\dim A$, and both $\mathcal V$ and $y_0$ pull
back to $\R^q$-valued maps on it.
On a precompact open set $\Omega$ over one half of $A$, the zeros of their
scaled difference $(\mathcal V-y_0)/R$ are precisely the preceding
intersections. A cutoff equal to one on $\Omega$ makes this difference a
smooth compactly supported map $G:Z\to\R^q$ satisfying
\[
 \deg(G,\Omega,0)\ne0,
 \qquad \|G\|_\infty>3\quad\text{on }\partial\Omega,
\]
where $\|\cdot\|_\infty$ is the maximum of the absolute target coordinates.
The scalar coordinates satisfy
\begin{equation}\label{eq:introduction-data}
 \max_i|d_{\widetilde M}G_i|\le J/R,
 \qquad \log q\le JR,
\end{equation}
with $J$ independent of $R$. Here $d_{\widetilde M}$ differentiates in the
fibre directions, holding the parameter in $A$ fixed, and its norm uses the
lifted metric of $M$. Parameter derivatives are bounded at each fixed $R$
and are controlled later by metric enlargement.
\Cref{sec:comparison} constructs these data without assuming that
$\widetilde M$ is contractible or that the cocycle is bounded.

\paragraph{A graph and one warped circle per coordinate.}
A graph change increases the fibre metric in the direction of one function's
differential, and a reciprocal warped circle supplies favorable scalar terms.
A nonlinear elliptic equation chooses a replacement coordinate close to the
old one while shortening its differential. Processing all $q$ coordinates
once is called a \emph{pass}. Paying an independent worst-case scalar cost
for every coordinate would introduce a factor $q$ and destroy the scalar
lower bound. Instead, the decreases of the inverse metric telescope on the
$n$-dimensional fibre cotangent space. Their sum has trace at most $n$ in the
metric at the start of the pass, independently of the numbers of coordinates
and circles. The scalar identity in \cref{sec:graph-calculus}, the graph
solution in \cref{sec:graph-existence}, and the summation in
\cref{sec:shortening} bound the fibre scalar-curvature loss by
$C(n+1)(2+\log q)L^2/w^2$, for an absolute constant $C$, when the input row
norms are at most $L$ and each replacement is allowed coordinate error
$\pi w/2$.

\paragraph{Finite shortening with a summable cost.}
Successive passes divide $L$ by four and $w$ by two, so their coordinate
errors and fibre scalar losses form summable series. The initial error
allowance makes the uniform total error in each coordinate less than $\pi/10$,
preserving the maximum-norm boundary gap. Composing the shortened coordinate
map on $\Omega$ with a degree-one map $\R^q\to S^q$ that is constant outside
the unit ball, and extending constantly outside $\Omega$, preserves this
degree. By \eqref{eq:introduction-data},
the total fibre scalar loss is $O((2+\log q)R^{-2})=O(R^{-1})$.
At each fixed finite stage, enlarging
the metric on $A$ makes the parameter part of the differential small and
makes the assembled scalar curvature uniformly close to the fibre scalar
curvature. For $R$ sufficiently large and every finite integer $k\ge1$,
there are a metric $g_k$ on $Z$, positive warps $m_{k,\nu}$, and a smooth map
$F_k:Z\to S^q$, constant outside a compact set and of nonzero degree, such
that
\[
 \Scal_{g_k+\sum_{\nu=1}^{kq}m_{k,\nu}^2d\theta_\nu^2}\ge\kappa/2,
 \qquad
 |dF_k|_{g_k,\mathrm{op}}\le C(q)4^{-k}J/R.
\]
The scalar bound concerns $Z\times\T^{kq}$, while the differential bound
uses all the base directions in $Z$.
\Cref{fig:direction-splitting} distinguishes these directions from the
stabilizing circles.

\begin{figure}[tb]
\centering
\begin{tikzpicture}[>=Latex,font=\small,
 block/.style={draw,rounded corners=2pt,minimum height=1.2cm,align=center}]
 \node[block,fill=blue!5,minimum width=3.2cm] (m) at (0,0)
   {Original fibre directions\\$\widetilde M$, dimension $n$};
 \node[block,fill=orange!7,minimum width=3.3cm] (a) at (3.65,0)
   {Parameter directions\\$A$, dimension $q-n$};
 \node[block,fill=purple!5,minimum width=3.0cm] (t) at (7.2,0)
   {Stabilizing circles\\$\T^{kq}$, dimension $kq$};
 \draw[thick] (-1.6,-.8)--(-1.6,-.95)--(5.3,-.95)--(5.3,-.8);
 \node[align=center] (z) at (1.85,-1.45)
   {Base $Z$, dimension $q$};
 \draw[->] (z.south)--node[right] {$F_k$, nonzero degree} (1.85,-2.65);
 \node at (1.85,-2.9) {$S^q$};
 \node[align=center] at (7.2,-1.6)
   {Scalar curvature is controlled\\on $Z\times\T^{kq}$};
\end{tikzpicture}
\caption{The local direction splitting after $k$ finite passes. The base
$Z$ is locally a product of original and parameter directions; no global
product decomposition is asserted. The inverse-metric trace uses only the
$n$ original directions. The sphere map uses the full base of dimension
$q$, which is fixed before the number $kq$ of circles is chosen.}
\label{fig:direction-splitting}
\end{figure}


\paragraph{An obstruction independent of the circle count.}
Pulling back a torus tube in $S^q$ gives a compact band $W\subset Z$ carrying
the degree. Doubling $W$ and retaining the circle factors gives a closed
oriented manifold $\operatorname{dbl}(W)\times\T^{kq}$ mapped with nonzero
degree to $S^1\times\T^{q-1}\times\T^{kq}$. The local torical theorem of
Cecchini--Schick requires a
scalar lower bound and a bound on the full differential only on the full
preimage of one connected open target arc. The doubled metric therefore needs
to agree with the constructed metric only there. Constant enlargement of the
stabilizing circles controls their joint projection without scalar loss.
\Cref{sec:obstruction} derives
an obstruction threshold depending on $q$ and $\kappa$, but not on the
number $kq$ of circles. We first choose $R$ large enough to control the total
fibre scalar loss, thereby fixing $q$ and all comparison data, and then choose a
sufficiently large finite $k$ to contradict this obstruction.

\subsection{Consequences}

For a connected aspherical manifold, the classifying map is a homotopy
equivalence. Thus every closed connected oriented aspherical manifold is
rationally essential. The main theorem gives the following consequence, with the
nonorientable case obtained from the orientation double cover.

\begin{corollary}[The Gromov--Lawson aspherical conjecture]
\label{cor:aspherical}
No closed smooth aspherical manifold, whether orientable or not, admits a
Riemannian metric of strictly positive scalar curvature.
\end{corollary}

For nonnegative scalar curvature, the corresponding statement determines
the given metric.

\begin{corollary}[Aspherical nonnegative-scalar-curvature rigidity]
\label{cor:aspherical-flat}
Let $g$ be a Riemannian metric of nonnegative scalar curvature on a closed
smooth aspherical manifold $M$, not necessarily oriented. Then $g$ is flat.
\end{corollary}

Both aspherical statements apply componentwise and have no spin assumption.
In dimensions zero and one scalar curvature vanishes
identically, and every metric is flat. Gromov already noted the rigidity
implication conditionally on the aspherical positive-scalar-curvature
obstruction \cite[Section~3, following Conjecture~B]{GromovDozen2018}.
Here the proof combines \cref{lem:scalar-reduction} with the compact-manifold
consequence of the Cheeger--Gromoll splitting theorem.

A nonzero-degree map to an aspherical target forces rational essentiality by
naturality of classifying maps; no injectivity or surjectivity condition on
the induced map of fundamental groups is needed. The following is the
oriented form of the target-map conjecture recorded by Chodosh and Li
\cite[Conjecture~18]{ChodoshLiSurvey2023}. Their survey also states the
connected-sum obstruction in dimensions four and five as a consequence of
Gromov's work \cite[Theorem~17]{ChodoshLiSurvey2023}.

\begin{corollary}[Oriented target-map obstruction]
\label{cor:degree-aspherical}
Let $M^n$ and $N^n$ be closed connected oriented smooth manifolds of
dimension $n\ge2$, with $N$ aspherical. If a continuous map $f:M\to N$
has nonzero ordinary oriented degree $d\in\Z\setminus\{0\}$, then $M$
admits no Riemannian metric whose scalar curvature is strictly positive
everywhere.

In particular, if $P^n$ is any closed connected oriented smooth manifold,
the oriented connected sum $N\#P$ admits no such metric.
\end{corollary}

For a closed connected oriented $n$-manifold, its \emph{real simplicial
volume} is
\[
 \|M\|=\inf\left\{\sum_i|a_i|:
       \sum_i a_i\sigma_i\text{ is a real singular cycle representing }[M]\right\}.
\]
Positive simplicial volume implies rational essentiality, by the
bounded-cohomology mapping theorem and duality for this seminorm
(\cref{lem:essentiality}). Therefore rational inessentiality implies zero
simplicial volume. The converse fails: a positive-dimensional torus is
rationally essential but has zero simplicial volume, as follows from its
self-covers of degree greater than one. Its flat metric also shows why
strict positivity in \cref{thm:main} cannot be replaced by nonnegative
scalar curvature. For simplicial volume the nonnegative hypothesis does
suffice.

\begin{corollary}[Simplicial-volume vanishing]
\label{cor:simplicial}
Every closed oriented smooth manifold of positive dimension
admitting a Riemannian metric of nonnegative scalar curvature has zero
real simplicial volume.
\end{corollary}

On each connected component, \cref{sec:reductions} reduces a
nonnegative-scalar-curvature metric to either a positive-scalar-curvature
metric or a Ricci-flat metric. The positive branch uses \cref{thm:main};
the Ricci-flat branch uses volume comparison and amenability, without
requiring rational inessentiality in that case.

\subsection{The obstruction problem and earlier work}

Two approaches have shaped scalar-curvature obstructions. Schoen and Yau
used stable two-sided minimal hypersurfaces and conformal deformation to
transfer positive scalar curvature to lower dimensions
\cite{SchoenYauStructure1979,SchoenYauIncompressible1979}. Their original
obstructions requiring no spin assumption apply in dimensions at most seven,
including to manifolds mapping with nonzero degree to tori. Gromov and Lawson
developed Dirac-operator obstructions using enlargeability and relative
indices on complete manifolds \cite{GromovLawson1983}. Their 1983
enlargeability convention requires oriented spin coverings, and their compact
obstruction has no upper dimension bound. The later geometric convention
used above omits this spin requirement. The spin condition in geometric
index applications and the regularity of minimizing hypersurfaces are thus
different constraints.

More recently, Chodosh and Li established the aspherical obstruction in
dimensions four and five using generalized soap bubbles
\cite{ChodoshLi2024}. Gromov independently treated the closed
five-dimensional case, and also complete-manifold domination problems under
properness and uniform-positivity or decay hypotheses
\cite{GromovAsphericalFive2020}.
\Cref{cor:aspherical} establishes the closed-aspherical statement in every
dimension.

The homological formulation in \cref{thm:main} concerns more manifolds than
the aspherical class. Hanke explicitly formulates the conjecture that
rationally essential manifolds do not admit positive scalar curvature
\cite[Definition~2.8 and Conjecture~2.9]{Hanke2011}. It is also the
classifying-map form, for closed oriented sources, of Gromov's conjecture
that a nontorsion homology class in the classifying space of a countable group
cannot be the image of the fundamental class of a positive-scalar-curvature manifold
\cite[Section~1.5, \(\lbrack\mathrm{Sc}\not>0\rbrack\), p.~19]{Gromov2023}.
The numbered sections cited from that work refer throughout to the author's
arXiv version~6. Yau also posed this homological obstruction in his survey
\cite[p.~265]{Yau2000}.

Gromov introduced simplicial volume and its bounded-cohomology description
in \cite{Gromov1982}. His quantitative scalar-curvature/simplicial-volume
conjecture asks for a dimension-dependent upper bound on simplicial volume in
terms of Riemannian volume under a fixed scalar-curvature lower bound
\cite[Section~3.13]{Gromov2023}. Rescaling a nonnegative-scalar-curvature
metric in such a comparison would force zero simplicial volume.
\Cref{cor:simplicial} proves this qualitative vanishing consequence.
The quantitative comparison remains a separate question.

Recent work of Ma, Wang, Xie, Yu and Zhu proves that nonnegative scalar
curvature forces zero real simplicial volume for closed oriented manifolds
with spin universal cover \cite[Corollary~1.2]{MaWangXieYuZhu2026}. Their
stated version has no further group-decay assumption. This conclusion
controls the simplicial norm, rather than the rational classifying image.
Bi and Zhu develop singular hypersurface descent and rule out positive scalar
curvature on closed oriented enlargeable and overtorical manifolds, where an
$n$-manifold is overtorical if it admits a nonzero-degree map to $\T^n$
\cite[Theorems~A and~B]{BiZhu2026}. Their theorem for complete oriented
sources and complete oriented enlargeable targets retains properness,
nonzero-degree and Lipschitz hypotheses and rules out a uniform positive
scalar lower bound.
Under strict positive scalar curvature,
\cref{thm:main} gives rational classifying-image vanishing without these
geometric-largeness or spin-cover hypotheses.

The terminal geometric input here is the local torical theorem of Cecchini
and Schick \cite[Theorem~1.14(1)]{CecchiniSchick2021}. It has no spin assumption
or upper dimension bound, and requires both the scalar lower bound and the
full differential bound only on the full preimage of one connected open target
arc. Its higher-dimensional proof uses the minimal-slicing theory of Schoen
and Yau \cite{SchoenYau2017}.

\Cref{sec:reductions} records the degree conventions and the nonnegative
scalar-curvature reduction. \Cref{sec:comparison} constructs the comparison
data, and \cref{sec:graph-calculus,sec:graph-existence,sec:shortening}
establish the graph stabilization and shortening estimates.
\Cref{sec:obstruction} supplies the stabilized obstruction.
Finally, \cref{sec:main-proof} combines these results and proves the four
corollaries.
